\documentclass{article}
\usepackage{graphicx} % Required for inserting images
\usepackage[utf8]{inputenc}
\usepackage[T1]{fontenc}
\usepackage[french]{babel}
\usepackage{amsmath,amssymb}
\usepackage{geometry}
\geometry{margin=2.5cm}

\title{Projet  RO : "transports de produits chimiques" }
\author{Jaber Tanouti \\ Rayane Rekik \\ Reda Atonane }
\date{2025 - 2026} 

\begin{document}

\maketitle

\section{Introduction}

Ce rapport étudie la gestion d'une flotte de camions pour transporter des
produits chimiques sur un horizon de cinq ans. Les acides partent de Liège vers
plusieurs villes belges, tandis que les bases sont transportées d'Anvers vers
Liège.

Le problème est décomposé en deux niveaux. Au niveau opérationnel, on fixe des
tournées-types réalistes et on en déduit les capacités utilisables. Au niveau
stratégique, on formule un modèle de flotte agrégé qui décide des achats, ventes
et affectations de camions. Le scénario opérationnel retenu pour paramétrer le
modèle principal est le scénario hybride S3, basé sur un couplage partiel
acide/base via Anvers.

\section{Indices et ensembles}

\begin{itemize}
  \item $t \in T = \{1,2,3,4,5\}$ : années de l'horizon de planification.
  \item $k \in K = \{1,2\}$ : types de camions :
  \begin{itemize}
    \item $k=1$ : camion type 1,
    \item $k=2$ : camion type 2.
  \end{itemize}
  \item $i,j \in V = \{\mathrm{AN},\mathrm{LI},\mathrm{CH},\mathrm{GA},\mathrm{BR},\mathrm{HA}\}$ : villes du réseau :
  \begin{itemize}
    \item LI = Liège, AN = Anvers, CH = Charleroi, GA = Gand, BR = Bruxelles, HA = Hasselt.
  \end{itemize}
  \item $J = V \setminus \{\mathrm{LI}\}$ : villes clientes pour les acides.
\end{itemize}

\section{Paramètres (constantes)}

\subsection{Données géographiques et temps}

\begin{itemize}
  \item $d_{i,j}$ : distance (km) entre deux villes $i,j \in V$. Les distances
  sont supposées symétriques. Les valeurs utiles de l'énoncé sont :
  \[
  \begin{array}{c|cccccc}
      & \mathrm{AN} & \mathrm{CH} & \mathrm{LI} & \mathrm{GA} & \mathrm{BR} & \mathrm{HA} \\
    \hline
    \mathrm{AN} & 0 & 100 & 105 & 40 & 45 & 50 \\
    \mathrm{CH} & 100 & 0 & 100 & 100 & 60 & 80 \\
    \mathrm{LI} & 105 & 100 & 0 & 140 & 100 & 60 \\
    \mathrm{GA} & 40 & 100 & 140 & 0 & 40 & 60 \\
    \mathrm{BR} & 45 & 60 & 100 & 40 & 0 & 50 \\
    \mathrm{HA} & 50 & 80 & 60 & 60 & 50 & 0
  \end{array}
  \]
  \item $v = 70$ km/h : vitesse moyenne des camions.
  \item $t^{\text{stop}} = 1$ h : durée d'un arrêt de livraison.
  \item $H^{\max} = 2\,000$ h/an : nombre maximal d'heures de travail par
  camion, avec 250 jours ouvrables et 8 h/jour.
  \item $\tau^{\text{change}} = 3$ jours ouvrables : durée d'immobilisation pour changement d'affectation d'un compartiment.
\end{itemize}

Dans le modèle principal, l'affectation des compartiments est fixée à l'année.
On ne déclenche donc pas de changement d'affectation en cours d'année et les
3 jours d'immobilisation ne sont pas retranchés des tournées journalières.

\subsection{Capacités physiques et contrainte légale}

\begin{itemize}
  \item Camion type 1 : un compartiment
  \[
    Q_1 = 16{,}5~\text{t}.
  \]
  \item Camion type 2 : deux compartiments
  \[
    Q_{2,16.5} = 16{,}5~\text{t} \quad\text{(grand compartiment)}, \qquad
    Q_{2,5.5} = 5{,}5~\text{t} \quad\text{(petit compartiment)}.
  \]
  \item Contrainte l\'egale : quantité maximale d'un m\^eme produit par camion
  \[
    Q^{\max} = 16{,}5~\text{t}.
  \]
  \item Quantité minimale par livraison :
  \[
    \underline{q}^{\text{liv}} = 5~\text{t}.
  \]
\end{itemize}

\subsection{Flotte initiale}

\[
  N_{1,0} = 4, \qquad N_{2,0} = 6.
\]

\subsection{Demandes annuelles}

\paragraph{Bases (Anvers $\rightarrow$ Liège).}

Pour tout $t \in T$ :
\[
  D^{B}_t = 30\,000 \quad \text{(t/an)}.
\]

\paragraph{Acides (Liège $\rightarrow$ autres villes).}

Pour tout $t \in T$ :
\begin{align*}
  D^{A}_{\mathrm{AN},t} &= 9\,000, &
  D^{A}_{\mathrm{CH},t} &= 12\,000, \\
  D^{A}_{\mathrm{GA},t} &= 2\,000, &
  D^{A}_{\mathrm{BR},t} &= 6\,200.
\end{align*}

Pour Hasselt :
\[
  D^{A}_{\mathrm{HA},1} = 350, \qquad
  D^{A}_{\mathrm{HA},2} = 825, \qquad
  D^{A}_{\mathrm{HA},t} = 1\,300 \quad \text{pour } t=3,4,5.
\]
La valeur de l'année 2 correspond à une demi-année à 350 t/an et une demi-année
à 1 300 t/an, car la nouvelle unité de Hasselt démarre après 18 mois.

Total acide à transporter en année $t$ :
\[
  \text{TotalAcide}_t = \sum_{j \in J} D^{A}_{j,t}.
\]

\subsection{Capacités annuelles moyennes (niveau 1)}

Ces paramètres intègrent la loi (16,5 t max/produit/camion), les 5 t minimum par livraison,
les distances, la vitesse, $t^{\text{stop}}$ et $H^{\max}$. L'immobilisation
$\tau^{\text{change}}$ n'est pas intégrée aux tournées du modèle principal car
les affectations sont fixes à l'année.

\begin{itemize}
  \item $K^{A,16.5}$ : tonnes d'acide/an livrables par un compartiment de 16{,}5 t.
  \item $K^{A,5.5}$ : tonnes d'acide/an livrables par un compartiment de 5{,}5 t.
  \item $K^{B,16.5}$ : tonnes de base/an livrables par un compartiment de 16{,}5 t.
  \item $K^{B,5.5}$ : tonnes de base/an livrables par un compartiment de 5{,}5 t.
\end{itemize}

Ces valeurs sont calculées en amont à partir des tournées-types du scénario S3.
Les scénarios S1 et S2 sont utilisés comme comparaisons, mais ne paramètrent pas
le modèle principal.

\subsection{Définition de la capacité annuelle $K^{A,16.5}$}

On considère un compartiment de capacité $16{,}5$ t dédié au transport d'acide.

\paragraph{1. Tournée type.}

On définit une tournée type au départ de Liège :
\[
  \text{LI} \to j_1 \to j_2 \to \dots \to j_m \to \text{LI},
\]
où
\[
  1 \le m \le 3, \quad j_\ell \in J, \; \ell = 1,\dots,m,
\]
et $m$ est le nombre de villes desservies lors de la tournée.

On note $q_\ell$ la quantité d'acide livrée à la ville $j_\ell$ (en tonnes).  
Les contraintes de quantité sont :
\begin{align}
  q_\ell &\ge \underline{q}^{\text{liv}} \quad &&\forall \ell = 1,\dots,m, \label{eq:q_min}\\
  \sum_{\ell=1}^{m} q_\ell &\le Q^{\max}. \label{eq:q_max}
\end{align}

\paragraph{2. Distance totale de la tournée.}

La distance totale parcourue lors de la tournée est :
\begin{equation}
  D^{A,\text{tour}}
  =
  d_{\mathrm{LI},j_1}
  + \sum_{\ell=1}^{m-1} d_{j_\ell, j_{\ell+1}}
  + d_{j_m,\mathrm{LI}}.
\end{equation}

\paragraph{3. Temps de trajet et temps d'arrêt.}

Le temps de trajet (sans les arrêts) est :
\begin{equation}
  T^{\text{trajet}} = \frac{D^{A,\text{tour}}}{v}.
\end{equation}

Le temps d'arrêt total (un arrêt par livraison) est :
\begin{equation}
  T^{\text{arrêts}} = m \cdot t^{\text{stop}}.
\end{equation}

\paragraph{4. Temps total par tournée.}

Le temps total consommé par une tournée type est donc :
\begin{equation}
  T^{A,\text{tour}}
  =
  T^{\text{trajet}} + T^{\text{arrêts}}
  =
  \frac{D^{A,\text{tour}}}{v} + m \, t^{\text{stop}}.
\end{equation}

\paragraph{5. Temps annuel disponible.}

On note $H^{\text{eff}}$ le temps annuel effectif (en heures/an) disponible pour ce compartiment :
\begin{equation}
  H^{\text{eff}} \le H^{\max},
\end{equation}
où $H^{\max}$ est le temps de travail maximal théorique (par exemple ajusté si l'on tient compte de jours d'immobilisation pour changement d'affectation).

\paragraph{6. Nombre de tournées par an.}

Le nombre maximal théorique de tournées de ce type que peut effectuer le compartiment est :
\begin{equation}
  N^{A,\text{tour/an}}
  =
  \frac{H^{\text{eff}}}{T^{A,\text{tour}}}
  =
  \frac{H^{\text{eff}}}{\frac{D^{A,\text{tour}}}{v} + m \, t^{\text{stop}}}.
\end{equation}

\paragraph{7. Quantité d'acide par tournée.}

La quantité totale d'acide livrée par une tournée est :
\begin{equation}
  Q^{A,\text{tour}} = \sum_{\ell=1}^{m} q_\ell,
\end{equation}


\paragraph{8. Capacité annuelle $K^{A,16.5}$.}

La capacité annuelle (en tonnes/an) d'un compartiment de 16{,}5 t dédié à l'acide est alors :
\begin{equation}
  K^{A,16.5}
  =
  Q^{A,\text{tour}} \times N^{A,\text{tour/an}}
  =
  Q^{A,\text{tour}} \cdot
  \frac{H^{\text{eff}}}{\frac{D^{A,\text{tour}}}{v} + m \, t^{\text{stop}}}.
\end{equation}

\subsection{Construction journalière progressive d'une flotte acide}

Pour dimensionner concrètement la flotte affectée aux acides, on commence par un
raisonnement journalier. On retient les demandes moyennes journalières des années
3 à 5, car elles correspondent au cas le plus contraignant avec Hasselt à pleine
demande. Avec $250$ jours ouvrables par an, on obtient :
\[
\begin{array}{c|ccccc|c}
\text{Ville} & \mathrm{AN} & \mathrm{CH} & \mathrm{GA} & \mathrm{BR} & \mathrm{HA} & \text{Total} \\
\hline
\text{Demande journalière (t/j)} & 36 & 48 & 8 & 24{,}8 & 5{,}2 & 122
\end{array}
\]

Un camion dispose de $8$ h par jour. Une tournée part de Liège, dessert une ou
plusieurs villes, puis revient à Liège. Le temps consommé par une tournée $r$ est :
\[
  T_r = \frac{D_r}{70} + n_r,
\]
où $D_r$ est la distance totale de la tournée et $n_r$ le nombre de villes livrées
pendant cette tournée. Chaque tournée transporte au maximum $16{,}5$ t d'acide.
Le modèle principal ne rajoute pas explicitement 0,5 h de rechargement à Liège
entre deux tournées ; cette variante est traitée comme analyse de sensibilité.

\paragraph{Cas d'un seul camion.}
La tournée aller-retour la plus courte est Liège--Hasselt--Liège :
\[
  T_{\mathrm{LI-HA-LI}} = \frac{120}{70} + 1 = 2{,}714 \text{ h}.
\]
Trois tournées même très courtes nécessiteraient $3 \times 2{,}714 = 8{,}142$ h,
donc un camion ne peut faire au maximum que deux tournées complètes par jour.
La quantité journalière maximale d'un camion est donc :
\[
  2 \times 16{,}5 = 33 \text{ t/j}.
\]
Un trajet optimal possible, en privilégiant les villes à forte demande, est :
\[
\mathrm{LI}\to\mathrm{BR}\to\mathrm{LI}
\quad\text{puis}\quad
\mathrm{LI}\to\mathrm{CH}\to\mathrm{LI}.
\]
Son temps total vaut :
\[
  \left(\frac{200}{70}+1\right) + \left(\frac{200}{70}+1\right)
  = 7{,}714 \text{ h},
\]
et il livre $33$ t. Un seul camion ne suffit donc pas pour couvrir les
$122$ t/j demandées.

\paragraph{Ajout progressif de camions.}
On ajoute ensuite les camions un par un. Pour chaque nombre de camions, on cherche
la quantité maximale livrable dans la journée, sous les contraintes suivantes :
\[
  \sum_{r \in R_c} T_r \le 8 \quad \forall c,
  \qquad
  \sum_{j \in r} q_{r,j} \le 16{,}5 \quad \forall r,
  \qquad
  0 \le \sum_r q_{r,j} \le D^{\text{jour}}_j.
\]

\[
\begin{array}{c|c|c|l}
\text{Nombre de camions} & \text{Quantité max. livrée (t/j)}
& \text{Demande restante (t/j)} & \text{Conclusion} \\
\hline
1 & 33 & 89 & \text{insuffisant} \\
2 & 66 & 56 & \text{insuffisant} \\
3 & 97{,}5 & 24{,}5 & \text{insuffisant} \\
4 & 114 & 8 & \text{encore insuffisant} \\
5 & 122 & 0 & \text{demande journalière couverte}
\end{array}
\]

Le seuil théorique basé uniquement sur la capacité serait
$\lceil 122/33\rceil = 4$ camions. Cependant, quatre camions ne suffisent pas en
pratique, car les demandes doivent être livrées dans des villes différentes et
certaines petites demandes, comme Gand et Hasselt, consomment du temps de tournée
sans remplir entièrement un camion.

\paragraph{Planning journalier retenu avec cinq camions.}
Le plan suivant couvre exactement la demande journalière moyenne des années 3 à 5.
Il respecte la limite de $8$ h par camion et garde chaque livraison physique au
moins égale à $5$ t.

\[
\begin{array}{c|l|c|c}
\text{Camion} & \text{Tournées et quantités livrées} & \text{Temps (h)} & \text{Total (t)} \\
\hline
1 &
\mathrm{LI}\!-\!\mathrm{CH}\!-\!\mathrm{LI}:16{,}5
\quad ; \quad
\mathrm{LI}\!-\!\mathrm{CH}\!-\!\mathrm{LI}:16{,}5
& 7{,}714 & 33 \\
2 &
\mathrm{LI}\!-\!\mathrm{CH}\!-\!\mathrm{LI}:15
\quad ; \quad
\mathrm{LI}\!-\!\mathrm{BR}\!-\!\mathrm{LI}:16{,}3
& 7{,}714 & 31{,}3 \\
3 &
\mathrm{LI}\!-\!\mathrm{AN}\!-\!\mathrm{LI}:12{,}35
\quad ; \quad
\mathrm{LI}\!-\!\mathrm{AN}\!-\!\mathrm{LI}:12{,}35
& 8{,}000 & 24{,}7 \\
4 &
\mathrm{LI}\!-\!\mathrm{AN}\!-\!\mathrm{HA}\!-\!\mathrm{LI}:
\mathrm{AN}\ 11{,}3,\ \mathrm{HA}\ 5{,}2
& 5{,}071 & 16{,}5 \\
5 &
\mathrm{LI}\!-\!\mathrm{GA}\!-\!\mathrm{BR}\!-\!\mathrm{LI}:
\mathrm{GA}\ 8,\ \mathrm{BR}\ 8{,}5
& 6{,}000 & 16{,}5
\end{array}
\]

La vérification par ville est :
\[
\begin{array}{c|ccccc|c}
\text{Ville} & \mathrm{AN} & \mathrm{CH} & \mathrm{GA} & \mathrm{BR} & \mathrm{HA} & \text{Total} \\
\hline
\text{Livré (t/j)} & 36 & 48 & 8 & 24{,}8 & 5{,}2 & 122
\end{array}
\]

On retient donc, pour le dimensionnement journalier des acides au régime permanent,
un besoin minimal de cinq camions de capacité $16{,}5$ t affectés aux acides.
Ce résultat correspond à la composante acide du scénario S3, qui est le scénario
opérationnel retenu pour le modèle stratégique.


\subsection{Coûts}

\begin{itemize}
  \item Coûts d'achat :
  \[
    C^{\text{achat}}_1 = 140\,000~\text{€}, \qquad
    C^{\text{achat}}_2 = 200\,000~\text{€}.
  \]
  \item Coût d'entretien annuel (par camion) :
  \[
    C^{\text{entretien}} = 5\,000~\text{€}/\text{an}.
  \]
  \item Prix d'achat pour la revente :
  \[
    C_1 = 140\,000, \qquad C_2 = 200\,000.
  \]
  \item $\alpha = 0{,}25$ : taux d'amortissement retenu dans le modèle principal.
  \item Prix théorique de vente d'un camion de type $k$ d'âge $n$ ans :
  \[
    C^{\text{vente}}_{k,n} = \frac{C_k}{(1+\alpha)^n}.
  \]
  \item Prix moyen de revente utilisé dans le modèle principal :
  \[
    \bar C^{\text{vente}}_{k,t}
    =
    \frac{C_k}{(1+\alpha)^t}.
  \]
\end{itemize}
Dans le modèle principal, nous ne suivons pas les cohortes d'achat des camions.
Nous utilisons donc un prix moyen de revente par type et par année. L'année $t$
est interprétée comme l'âge moyen du camion vendu en fin d'année $t$.

Avec $\alpha = 0{,}25$, on obtient :
\[
  \bar C^{\text{vente}}_{1,t}
  =
  \frac{140\,000}{1{,}25^t},
  \qquad
  \bar C^{\text{vente}}_{2,t}
  =
  \frac{200\,000}{1{,}25^t}.
\]

\begin{center}
\begin{tabular}{c|cc}
\hline
Année $t$ & $\bar C^{\text{vente}}_{1,t}$ & $\bar C^{\text{vente}}_{2,t}$ \\
\hline
1 & 112\,000 & 160\,000 \\
2 & 89\,600 & 128\,000 \\
3 & 71\,680 & 102\,400 \\
4 & 57\,344 & 81\,920 \\
5 & 45\,875{,}20 & 65\,536 \\
\hline
\end{tabular}
\end{center}
Une analyse de sensibilité pourra être menée avec
$\alpha \in \{0{,}15;0{,}25;0{,}35\}$.

\section{Variables de décision}

\subsection{Flotte}

Pour $k \in \{1,2\}$, $t \in T$ :
\begin{itemize}
  \item $N_{k,t} \in \mathbb{Z}_+$ : nombre de camions de type $k$ disponibles pendant l'année $t$.
  \item $A_{k,t} \in \mathbb{Z}_+$ : nombre de camions de type $k$ achetés au début de l'année $t$.
  \item $V_{k,t} \in \mathbb{Z}_+$ : nombre de camions de type $k$ vendus en fin d'année $t$.
\end{itemize}

\subsection{Affectation des camions type 1}

Pour chaque année $t$ :
\begin{itemize}
  \item $X^{A}_{1,t} \in \mathbb{Z}_+$ : nombre de camions type 1 affectés aux acides.
  \item $X^{B}_{1,t} \in \mathbb{Z}_+$ : nombre de camions type 1 affectés aux bases.
\end{itemize}

\subsection{Affectation des camions type 2 (configurations)}

Pour chaque année $t$ :
\begin{itemize}
  \item $Y^{AB}_t \in \mathbb{Z}_+$ : type 2 avec grand compartiment acide (16{,}5 t) et petit compartiment base (5{,}5 t).
  \item $Y^{BA}_t \in \mathbb{Z}_+$ : type 2 avec grand compartiment base (16{,}5 t) et petit compartiment acide (5{,}5 t).
  \item $Y^{A0}_t \in \mathbb{Z}_+$ : type 2 avec grand compartiment acide, petit compartiment vide.
  \item $Y^{B0}_t \in \mathbb{Z}_+$ : type 2 avec grand compartiment base, petit compartiment vide.
\end{itemize}
Ces configurations représentent le couplage partiel retenu via Anvers. Le
modèle principal n'autorise pas un couplage libre de tous les produits sur tous
les trajets.

\section{Contraintes}

\subsection{Évolution de la flotte}

Les achats sont disponibles dès l'année $t$, tandis que les ventes ont lieu en
fin d'année $t$. Un camion vendu en année $t$ reste donc utilisable pendant
l'année $t$ et disparaît de la flotte disponible à partir de l'année $t+1$.

Pour $k \in \{1,2\}$ et $t = 1,\dots,5$ :
\[
  N_{k,t} = N_{k,t-1} - V_{k,t-1} + A_{k,t},
\]
avec $V_{k,0}=0$.

Conditions initiales :
\[
  N_{1,0} = 4, \qquad N_{2,0} = 6.
\]

On impose aussi que les ventes de fin d'année ne dépassent pas la flotte
disponible pendant cette année :
\[
  V_{k,t} \le N_{k,t}
  \qquad \forall k \in \{1,2\},\ \forall t \in T.
\]

\subsection{Cohérence affectation / flotte}

\paragraph{Type 1.} Pour chaque année $t$ :
\[
  X^{A}_{1,t} + X^{B}_{1,t} = N_{1,t}.
\]

\paragraph{Type 2.} Pour chaque année $t$ :
\[
  Y^{AB}_t + Y^{BA}_t + Y^{A0}_t + Y^{B0}_t = N_{2,t}.
\]

\subsection{Capacité annuelle pour les acides}

Capacité annuelle totale pour les acides en année $t$ :
\[
  \text{Cap}^A_t
  =
  K^{A,16.5}\, X^{A}_{1,t}
  + K^{A,16.5}\, (Y^{AB}_t + Y^{A0}_t)
  + K^{A,5.5}\, Y^{BA}_t.
\]

Contrainte de capacité acide (satisfaction de la demande) :
\[
  \text{TotalAcide}_t \;\le\; \text{Cap}^A_t
  \qquad \forall t \in T.
\]

\subsection{Capacité annuelle pour les bases}

Capacité annuelle totale pour les bases en année $t$ :
\[
  \text{Cap}^B_t
  =
  K^{B,16.5}\, X^{B}_{1,t}
  + K^{B,16.5}\, (Y^{BA}_t + Y^{B0}_t)
  + K^{B,5.5}\, Y^{AB}_t.
\]

Demande de base : $D^{B}_t = 30\,000$ t/an.  
Contrainte de capacité base :
\[
  D^{B}_t \;\le\; \text{Cap}^B_t
  \qquad \forall t \in T,
\]
c'est-à-dire
\[
  30\,000 \;\le\; \text{Cap}^B_t
  \qquad \forall t \in T.
\]

\subsection{Intégralité et non-négativité}

Pour tout $t \in T$ :
\[
  N_{k,t},\, A_{k,t},\, V_{k,t} \in \mathbb{Z}_+ \quad (k = 1,2),
\]
\[
  X^{A}_{1,t},\, X^{B}_{1,t} \in \mathbb{Z}_+,
\]
\[
  Y^{AB}_t,\, Y^{BA}_t,\, Y^{A0}_t,\, Y^{B0}_t \in \mathbb{Z}_+.
\]

\section{Fonction économique}

On minimise le coût total sur les 5 années : achats + entretien $-$ recettes de vente.

Paramètres :
\[
  C^{\text{achat}}_1 = 140\,000~\text{€}, \qquad
  C^{\text{achat}}_2 = 200\,000~\text{€},
\]
\[
  C^{\text{entretien}} = 5\,000~\text{€}/\text{an}.
\]

$\bar C^{\text{vente}}_{k,t}$ : prix moyen de vente d'un camion de type $k$ en
fin d'année $t$, défini par
\[
  \bar C^{\text{vente}}_{k,t}
  =
  \frac{C_k}{(1+\alpha)^t}.
\]
Ce paramètre remplace un suivi détaillé par cohortes. Dans le modèle principal,
$\alpha = 0{,}25$.

Variables concernées : $A_{k,t}$, $V_{k,t}$, $N_{k,t}$.

Fonction objectif :
\[
  \min Z
  =
  \underbrace{\sum_{t=1}^{5} \sum_{k=1}^{2}
    C^{\text{achat}}_k \, A_{k,t}}_{\text{coût d'achat}}
  +
  \underbrace{\sum_{t=1}^{5} \sum_{k=1}^{2}
    C^{\text{entretien}} \, N_{k,t}}_{\text{coût d'entretien}}
  -
  \underbrace{\sum_{t=1}^{5} \sum_{k=1}^{2}
    \bar C^{\text{vente}}_{k,t} \, V_{k,t}}_{\text{recettes de vente}}.
\]

\end{document}
